\section{Sequential algorithm}

\subsection{Problem definition}
We formalize the Forest Fire simulation as a probabilistic cellular automaton operating on a square two-dimensional spatial grid $\mathcal{G}$ of dimensions $N \times N$. Each cell $(i, j)$ within the grid, where $0 \le i < N$ and $0 \le j < N$, can exist in exactly one of three mutually exclusive states at any given discrete time step $t$: $s \in \{EMPTY, TREE, FIRE\}$. 

The dynamics of the system are governed by localized interactions and stochastic processes. Specifically, we define two input probability parameters:
\begin{itemize}
    \item $p \in [0, 1]$: The probability of a new tree growing in an empty cell.
    \item $l \in [0, 1]$: The probability of a tree spontaneously igniting due to a lightning strike.
\end{itemize}
The propagation of the fire relies on the standard Von Neumann neighborhood, meaning that the state of any cell $(i, j)$ at time $t+1$ depends solely on its own state and the states of its four orthogonal neighbors (North, South, East, West) at time $t$.

\subsection{Algorithm and complexity}
In our baseline sequential implementation, we simulate the evolution of the grid over $T$ total time steps. At each time step, we apply the transition rules uniformly across all cells. The standard rules dictate that a burning cell consumes its fuel and becomes empty; a tree ignites if it is adjacent to a burning cell or catches fire spontaneously with probability $l$; and an empty space grows a tree with probability $p$.

The pseudocode for this sequential process is detailed in Algorithm \ref{alg:sequential_forest}.

\begin{algorithm}[H]
\caption{Sequential Forest Fire Simulation}
\label{alg:sequential_forest}
\begin{algorithmic}[1]
\REQUIRE Grid dimension $N$, probabilities $p, l$, number of steps $T$
\ENSURE Evolved grid states over time
\FOR{$t = 1$ \TO $T$}
    \STATE Allocate an empty grid $\mathcal{G}_{next}$ of size $N \times N$
    \FOR{each cell $(i, j)$ in $\mathcal{G}$}
        \IF{$\mathcal{G}[i][j] == FIRE$}
            \STATE $\mathcal{G}_{next}[i][j] \leftarrow EMPTY$
        \ELSIF{$\mathcal{G}[i][j] == TREE$}
            \IF{any orthogonal neighbor is $FIRE$}
                \STATE $\mathcal{G}_{next}[i][j] \leftarrow FIRE$
            \ELSE
                \STATE $\mathcal{G}_{next}[i][j] \leftarrow FIRE$ with probability $l$
                \STATE $\mathcal{G}_{next}[i][j] \leftarrow TREE$ with probability $(1 - l)$
            \ENDIF
        \ELSIF{$\mathcal{G}[i][j] == EMPTY$}
            \STATE $\mathcal{G}_{next}[i][j] \leftarrow TREE$ with probability $p$
            \STATE $\mathcal{G}_{next}[i][j] \leftarrow EMPTY$ with probability $(1 - p)$
        \ENDIF
    \ENDFOR
    \STATE $\mathcal{G} \leftarrow \mathcal{G}_{next}$
\ENDFOR
\end{algorithmic}
\end{algorithm}

To ensure that the state transitions remain purely synchronous and to prevent data hazards (such as read-after-write conflicts during a single iteration), we employ a double-buffering technique. The new states are written to an auxiliary grid, $\mathcal{G}_{next}$, before being copied back to the primary grid $\mathcal{G}$ at the end of the time step.

\textbf{Complexity Analysis:} 
Because we must maintain two full grid structures simultaneously to preserve the strict data-dependency constraints, the space complexity of our sequential algorithm is strictly bounded by $\mathcal{O}(N^2)$. Regarding time complexity, calculating the next state for a single cell requires a constant $\mathcal{O}(1)$ number of operations (evaluating the four neighbors and generating pseudo-random numbers). Since we must traverse the entire $N \times N$ grid at each time step, the temporal complexity per iteration is $\mathcal{O}(N^2)$. Thus, the overall time complexity for the complete simulation over $T$ steps is $\mathcal{O}(T \times N^2)$.